% Generated draft: requires mathematical and source review.
\documentclass[11pt]{article}
\usepackage[margin=25mm]{geometry}
\usepackage{fontspec}
\setmainfont{Latin Modern Roman}
\usepackage{amsmath,amssymb,mathtools,unicode-math}
\setmathfont{Latin Modern Math}
\usepackage{enumitem,needspace,xurl,hyperref}
\setlength{\parindent}{0pt}
\setlength{\parskip}{5pt}
\setlength{\emergencystretch}{3em}
\newcommand{\N}{\mathbb{N}}
\newcommand{\Z}{\mathbb{Z}}
\newcommand{\Q}{\mathbb{Q}}
\newcommand{\R}{\mathbb{R}}
\newcommand{\C}{\mathbb{C}}
\newcommand{\F}{\mathbb{F}}
\newcommand{\A}{\mathbb{A}}
\newcommand{\PP}{\mathbb{P}}
\newcommand{\E}{\mathbb{E}}
\newcommand{\eps}{\varepsilon}
\newcommand{\dd}{\,\mathrm{d}}
\begin{document}
\section*{UnsolvedMath ID: 1161}
\textbf{Source problem code:} catalogue entry
\subsection{Definitions and mathematical statement}
The supplied historical question is: give Schubert's enumerative calculus a rigorous foundation. This is a foundational program rather than a single quantified assertion in the record. A rigorous foundation is intended to define the geometric incidence conditions and their intersection products, specify when enumerative counts are finite and independent of generic choices, and account for intersection multiplicities and degenerations under stated hypotheses.

For a concrete illustration of the objects involved, let $K$ be a field, let $V$ be an $n$-dimensional $K$-vector space, and let $\operatorname{Gr}(r,V)$ be the Grassmannian of $r$-dimensional linear subspaces, where $0\leq r\leq n$. A complete flag is a chain $0=F_0\subset F_1\subset\cdots\subset F_n=V$ with $\dim F_i=i$. For a partition $\lambda=(\lambda_1,\ldots,\lambda_r)$ satisfying $n-r\geq\lambda_1\geq\cdots\geq\lambda_r\geq0$, its Schubert incidence locus is
\[
\Omega_\lambda(F_\bullet)=\{W\in\operatorname{Gr}(r,V):
\dim(W\cap F_{n-r+i-\lambda_i})\geq i\text{ for }1\leq i\leq r\}.
\]
Modern intersection-theoretic formulations attach classes to such loci in a Chow ring, or in a suitable cohomology ring when that theory is available, and interpret appropriate top-degree intersection products as enumerative numbers subject to hypotheses on the field, incidence data, properness and multiplicity. This is an illustrative precise setting for the foundational request, not a replacement claim that the source restricted Hilbert's problem to one field, one family of varieties, or one enumerative formula. Flag varieties are also included in the supplied modern-foundation assessment.

The background identifies this as Hilbert's fifteenth problem from his 1900 list. Its example asks about lines in three-dimensional space meeting four given lines, but it does not specify the field, affine or projective setting, genericity, or multiplicity convention, so no definite numerical count is asserted here. The background mentions intersection theory, Chow rings, Gromov--Witten invariants and quantum cohomology, and describes continuing research concerning classical Schubert calculus in all dimensions.

The supplied literature triage, checked 2026-08-17, instead assesses the historical foundational request as fulfilled by modern intersection theory and the cohomology of Grassmannians and flag varieties, labels it solved, and states that no work remains for the stated foundational problem. It records no separate verified partial result. This dated assessment is retained alongside the registry's primary ``open'' label and the background's research-oriented wording; it does not turn unspecified modern research questions into a new open conjecture.
\subsection{Short English statement}
Hilbert's fifteenth problem asks for rigorous definitions and justifications for Schubert's geometric counting calculus, including incidence conditions and intersection multiplicities. The supplied 2026-08-17 assessment says that the historical foundational request has been fulfilled; the record does not specify a separate remaining theorem to prove.
\subsection{Sources}
\begin{itemize}
\item UnsolvedMath local registry, permanent ID 1161, source code catalogue entry, ``Hilbert's Fifteenth Problem''. The supplied background attributes the foundational question to Hilbert's 1900 list. Dataset attribution: Ulam AI and the UnsolvedMath Contributors; curation and metadata under CC BY 4.0.
\item \emph{Schubert calculus}, overview, \url{https://en.wikipedia.org/wiki/Schubert_calculus}, recorded as accessed 2026-08-17. The local triage labels this an authoritative secondary source and cites its description of the modern intersection-theoretic formulation.
\item No primary historical text, foundational monograph, or precise theorem citation is supplied in this record. The modern solved assessment is reproduced from the local triage rather than independently verified.
\end{itemize}
\end{document}
